\chapter{Conclusion}
\label{chap:conclusion}

\section{Bilan}

Les adaptateurs, l'intercepteur OAuth2, l'appairage des passagers et la retarification tournent chez les agences suivies. \texttt{rawCarrierMetadata} est toujours là. La table d'exceptions ERP aussi. Quand le transporteur ne renvoie plus d'offre après un 409, le parcours reste médiocre.

J'ai passé beaucoup plus de temps que prévu sur la taxe IZ, les fuseaux Aegean/Turkish et le deadlock InnoDB. Depuis, je prépare plus souvent une repro avant de modifier un contrôleur, et je pose la question au PO quand une règle métier est floue.

\section{RNCP 36137}

Pour le titre d'\textbf{Expert en Ingénierie Informatique}, ce projet m'a mis devant une architecture multi-services, du full-stack Node.js / Nuxt, des contraintes ADM, des tests automatisés et une équipe Paris / Tunisie. Je n'ai pas déroulé le référentiel point par point ; ce sont les sujets que j'ai réellement traités.

\section{Perspectives}

Deux sujets à suivre : \textbf{IATA ONE Order} et \textbf{OpenTelemetry} sur les appels transporteurs.
